\section{The HPS experiment and data samples}
\label{sec:detector-data}

HPS is a compact forward spectrometer located downstream of a thin tungsten target in
Hall B at Jefferson Lab. The detector is optimized for boosted $e^+e^-$ pairs produced
at small angles relative to the beam direction. Its prompt-resonance sensitivity is
driven by the Silicon Vertex Tracker (SVT), which measures charged-particle momenta and
vertices, and by the electromagnetic calorimeter (ECal), which provides the fast
trigger and track--cluster matching \cite{HPSExperiment2022,HPS2016PromptLong}.

The data samples in this analysis span two detector eras. The 2015 and 2016
engineering runs used the baseline six-station SVT, with the first active silicon
station near the target and close to the beam plane. Before the 2019 and 2021 physics
runs, HPS installed an upgraded front tracker with an additional double layer,
slim-edge sensors, and reduced material in the innermost tracking stations
\cite{HPSExperiment2022,HPSSlimEdgeSensors2020}. The upgraded configuration improves
low-mass acceptance and changes the mass-resolution model used by the signal template.
The ECal layout remained the same basic subsystem, with a positron hodoscope added
before the physics runs to support upgraded triggering.

\begin{figure*}[t]
\centering
\begin{minipage}{0.48\textwidth}
\centering
\ifigure{\linewidth}{hps_svt_2015_2016.png}\\
(a) Baseline engineering-run SVT.
\end{minipage}
\hfill
\begin{minipage}{0.48\textwidth}
\centering
\ifigure{\linewidth}{hps_svt_2019_2021.png}\\
(b) Upgraded 2019/2021 front tracker.
\end{minipage}
\caption{Detector-era comparison relevant to the prompt-resonance analysis. The 2015
and 2016 samples use the baseline six-station SVT, while the 2021 validation sample
uses the upgraded front tracker. The detector change is a direct input to the
mass-resolution and acceptance model used by the GPR workflow.}
\label{fig:svt-eras}
\end{figure*}

Table~\ref{tab:dataset-summary} summarizes the samples carried through the current
workflow. The 2015 sample is fully unblinded and is used as the most mature validation
anchor. The 2016 sample is represented here by the 10\% development subset of the
full 2016 parent sample. The 2021 sample is represented by the 1\% validation subset
of the upgraded-detector physics run. These partial samples are useful for validating
the combined method, but they are not the final full-statistics result.

\begin{table*}[t]
\caption{\label{tab:dataset-summary}Data samples and current GPR scan ranges. The
quoted 2016 and 2021 fractions describe the staged validation samples used in the
current workflow, not the final full-statistics publication data set.}
\small
\begin{tabular*}{\textwidth}{@{\extracolsep{\fill}}p{0.18\textwidth}p{0.11\textwidth}p{0.18\textwidth}p{0.13\textwidth}p{0.28\textwidth}@{}}
\toprule
Sample & Beam energy & Exposure represented here & GPR scan range & Role in this draft \\
\midrule
2015 engineering run & 1.056 GeV & 1170 nb$^{-1}$ & 20--130 MeV & Fully unblinded validation anchor \\
2016 engineering run & 2.3 GeV & 10\% of 10608 nb$^{-1}$ & 35--210 MeV & Development subset of the 2016 prompt sample \\
2021 physics run & 3.74 GeV & 1\% of $\sim$160 pb$^{-1}$ & 30--250 MeV & Upgraded-detector validation subset \\
\bottomrule
\end{tabular*}
\end{table*}


\begin{figure*}[t]
\centering
\ifigure{0.86\textwidth}{input_mass_spectra.png}
\caption{Selected invariant-mass spectra used by the current GPR workflow. The spectra
are shown in absolute events per MeV for the 2015, 2016 10\%, and 2021 1\% samples.
The shaded regions mark the published HPS engineering-run search windows and the
current GPR scan regions.}
\label{fig:input-spectra}
\end{figure*}

At a common mass hypothesis, the active data samples are statistically independent.
They differ in beam energy, exposure, detector resolution, event selection, and
radiative-trident normalization. The combined analysis therefore treats them as
separate channels linked only by the same dark-photon coupling parameter $\eps^2$.
